\section{Introduction}
\label{sec:introduction}

The TPC-32 matched shell reduces one high-row fixed-shift problem to
the distinguished zero-frequency concentration of a literal
post-bin coefficient \(A_X\) \citep{WangTPC32}.  Writing
\[
 Z_X=\sum_nA_X(n),\qquad
 D_X=\sum_n|A_X(n)|^2,
\]
the statistic is
\[
 \Flat_X=\frac{|Z_X|^2}{D_X}.
\]
TPC-81 and TPC-82 separate the two arithmetic obligations behind
this quotient: determinant energy must survive, and the signed zero
mode must cancel \citep{WangTPC81,WangTPC82}.  TPC-88 then factors
the same statistic into post-bin magnitude spreading and relative
\(\ell^1\) phase cancellation \citep{WangTPC88}.

The present paper asks a different, operational question:

\begin{quote}
If the literal coefficient is replaced by an analytically or
numerically tractable surrogate, how accurately must that surrogate
approximate the coefficient before its zero-mode geometry transfers?
\end{quote}

There is a canonical Hilbert-space formulation.  On the actual
support \(\Omega\), let
\[
 e=\frac{\one}{\sqrt{\#\Omega}}.
\]
For a nonzero coefficient \(a\), define
\[
 \rho(a)=\frac{|\langle a,e\rangle|}{\|a\|_2}.
\]
This is the cosine of the projective angle from \(a\) to the constant
line.  Since
\[
 \langle a,e\rangle=\frac{1}{\sqrt{\#\Omega}}\sum_{\omega\in\Omega}
 a_\omega,
\]
one has the exact identity
\[
 \Flat(a)=\#\Omega\,\rho(a)^2.
\]
Thus the zero-frequency gate is a principal-angle gate.

The principal angle is especially well suited to surrogate
analysis because it is invariant under nonzero scalar
renormalization.  More importantly, its perturbation law is exact.
If \(a\) lies in the relative ball
\[
 \|a-b\|_2\le\delta\|b\|_2,\qquad \delta<1,
\]
then the projective angle between \(a\) and \(b\) is at most
\(\arcsin\delta\).  This yields sharp upper and lower envelopes for
\(\rho(a)\), and in particular
\[
 |\rho(a)-\rho(b)|\le\delta.
\]
The constant is best possible.

This elementary fact has a nontrivial consequence for the TPC
program.  The critical zero-mode scale is \(1/J_X\), not merely a
quantity tending to zero.  Therefore a surrogate with relative
\(\ell^2\) error \(o(1)\) may preserve its energy scale while
completely failing to preserve the required zero-mode scale.
Uniform transfer requires \(O(1/J_X)\) accuracy.  A two-dimensional
tangent construction shows that this order cannot be improved.

\subsection{Main contributions}

The paper establishes the following.
\begin{enumerate}[label=(\roman*),leftmargin=2.3em]
\item It identifies determinant flatness with the squared
principal-angle cosine on the prescribed actual support:
\[
 D=\frac{|Z|^2}{S}+D_\perp,\qquad
 \Flat=S\rho^2.
\]
\item It gives the sharp two-sided perturbation envelope for
\(\rho\) under relative \(\ell^2\) error, including all boundary
cases and equality constructions.
\item It audits the simpler quotient estimate
\[
 \rho(a)\le\frac{\rho(b)+\delta}{1-\delta}
\]
and shows that the exact projective result is stronger.
\item It proves sharp transfer inequalities for determinant
flatness, coherent energy, orthogonal energy, and the zero mode.
\item It proves that \(O(1/J)\) relative accuracy is the correct
uniform threshold for transferring \(\rho\ll1/J\), whereas \(o(1)\)
is not enough.
\item It shows that preserving a two-sided
\(\rho\asymp1/J\) statement needs a constant margin; an error of the
same size as the lower signal can erase that signal exactly.
\item It gives a literal same-support \(\Lone\) certificate for the
fixed-\(h_0\) TPC-32 coefficient.
\item It isolates zero padding as a normalization trap:
\(\rho\) changes under padding although \(\Flat\) does not.
\item It keeps the determinant-energy and physical-loss ledgers
independent.
\end{enumerate}

\subsection{Audit levels}

We use the established terminology.  An \(\Lzero\) result is finite
Hilbert-space algebra or an abstract finite counterexample.  An
\(\Lone\) result places that algebra on the literal fixed-\(h_0\)
coefficient with its actual bins and original normalization.  An
\(\Ltwo\) result is a new growing-\(X\) arithmetic estimate for that
coefficient.  All angle and perturbation theorems below are
\(\Lzero\).  The same-frame implication for \(A_X\) is \(\Lone\).
No surrogate satisfying its hypotheses is constructed, and no
\(\Ltwo\) estimate is claimed.
